\documentclass[11pt]{article}
\usepackage[T1]{fontenc}
\usepackage[margin=1in]{geometry}
\newcounter{exercise}
\title{Problem Set with an Appendix and a Custom Counter}
\author{FlashTeX Fixtures Team}
\date{}
\begin{document}
\maketitle
\tableofcontents
\section{Warm-up}
The starred challenge is Exercise~\ref{ex:star} on page~\pageref{ex:star}; attempt it only after the warm-ups. \refstepcounter{exercise}Exercise~\theexercise: differentiate $f(x) = x^3 \sin x$. \refstepcounter{exercise}Exercise~\theexercise: evaluate $\int_0^1 x e^x \, dx$. \setcounter{exercise}{10} The numbering now jumps: \refstepcounter{exercise}Exercise~\theexercise: solve $y' = y$ with $y(0) = 2$. \addtocounter{exercise}{4} And jumps again: \refstepcounter{exercise}Exercise~\theexercise: prove that $\sqrt{2}$ is irrational.
\section{Main problems}
\label{sec:main}
The appendix collects solutions, including the solution to the starred Exercise~\ref{ex:star}. \refstepcounter{exercise}\label{ex:star}Exercise~\theexercise: (starred) prove that there are infinitely many primes. The techniques of Section~\ref{sec:main} suffice for every problem above, and the solutions in Appendix~\ref{app:sol} on page~\pageref{app:sol} follow the same numbering.
\appendix
\section{Solutions}
\label{app:sol}
Solutions follow the exercise numbers fixed in the main text: the warm-ups, the jumped numbers produced by \texttt{setcounter} and \texttt{addtocounter}, and finally the starred Exercise~\ref{ex:star} of Section~\ref{sec:main}, whose proof is Euclid's classic argument.
\end{document}
